% Section 3. Every factual sentence ends with a comment naming its ledger rows (results/2026-10-26-claims-ledger/ledger-v7.tsv).
% Quotes: a line "%% QUOTE <path> <a>-<b> [sha=<ledger row>]" precedes each quoted block; check.py compares the block with those source lines.
\section{Statement faithfulness}\label{sec:faithfulness}

We read four Comparator challenge statements of family 003 against the paper statements they encode: SiegelZeros, QuasiRiemannHypothesis, DirichletSevenEighths and HeckeSevenEighths. % v7:L001 v7:L002 v7:L004 v7:L006 v7:L008
Each finding rests on reading the Lean statement against the paper's text, so each is INFERRED. % v7:L002 v7:L004 v7:L006 v7:L008
Table~\ref{tab:faithfulness} gives each challenge file's sha256 at openai/math fd4aeeb, the finding, and the record of the reading. % v7:L001 v7:L002 v7:L003 v7:L004 v7:L005 v7:L006 v7:L007 v7:L008

\subsection{The SiegelZeros challenge}

The SiegelZeros challenge file has sha256 bedaac0248a641910f79a380d9702d82c6bfbbfd3760a69d27c72c5d8fb136ce, and it is unchanged between adc7f12 and fd4aeeb. % v7:L007
Listing~\ref{lst:siegel-challenge} quotes its two theorems. % v7:L007
By reading the Lean statement against the paper's text, the challenge statement matches the paper's Theorem 1, with the hypotheses primitive, nonprincipal, real, q >= 3 and beta in (0,1). % v7:L008
The type of each solution theorem equals the type of the challenge theorem. % v7:L008

\begin{listing}[tbp]
\caption{The SiegelZeros challenge, lean/ComparatorChallenges/SiegelZeros.lean:8-14 and 19-25 (openai/math fd4aeeb; file sha256 bedaac02\ldots, unchanged since adc7f12).}\label{lst:siegel-challenge} % v7:L007
%% QUOTE external/openai-math/lean/ComparatorChallenges/SiegelZeros.lean 8-14 sha=v7:L007
\begin{Verbatim}
theorem exists_absolute_real_zero_gap :
    ∃ c : ℝ, 0 < c ∧
      ∀ (q : ℕ) [NeZero q], 3 ≤ q →
      ∀ χ : DirichletCharacter ℂ q,
        χ.IsPrimitive → χ ≠ 1 → (∀ a : ZMod q, (χ a).im = 0) →
        ∀ β : ℝ, 0 < β → β < 1 → χ.LFunction (β : ℂ) = 0 →
          c ≤ (1 - β) * Real.log (q : ℝ) := by
\end{Verbatim}
%% QUOTE external/openai-math/lean/ComparatorChallenges/SiegelZeros.lean 19-25 sha=v7:L007
\begin{Verbatim}
theorem WeightedTorusJets.dirichletRealZeroBound_proof :
    ∃ c : ℝ, 0 < c ∧ ∀ (q : ℕ) [NeZero q], 3 ≤ q →
      ∀ χ : DirichletCharacter ℂ q,
        χ.IsPrimitive → χ ≠ 1 → (∀ a : ZMod q, (χ a).im = 0) →
        ∀ β : ℝ,
          (0 < β ∧ β < 1 ∧ DirichletCharacter.LFunction χ (β : ℂ) = 0) →
            c ≤ (1 - β) * Real.log (q : ℝ) := by
\end{Verbatim}
\end{listing}

The paper states its Theorem 1 at refs/siegel-paper.tex:66-73. % v7:L106
Its text, lines 67-72, reads: % v7:L106
%% QUOTE refs/siegel-paper.tex 67-72 sha=v7:L143
\begin{paperquote}{1.1}
There is an absolute constant $c>0$ such that every real zero
$\beta\in(0,1)$ of every primitive nonprincipal real Dirichlet
$L$-function of conductor $q\ge3$ satisfies
\begin{equation}\label{eq:main}
 (1-\beta)\log q\ge c.
\end{equation}
\end{paperquote}

Four Lean lemmas show that the hypotheses of the challenge can be met and what the statement specialises to. % v7:L011 v7:L012 v7:L013 v7:L014
NonVacuity.hypotheses\_satisfiable proves that for every odd prime $p$ the quadratic character mod $p$ satisfies all hypotheses of the challenge ($q = p \ge 3$, primitive, nonprincipal, real). % v7:L011
NonVacuity.quadChar\_LFunction\_eq\_LSeries proves that the L-function of the quadratic character equals its Dirichlet series for $\operatorname{Re} s > 1$. % v7:L012
NonVacuity.quadChar\_LFunction\_one\_ne\_zero proves $L(1, \chi_p) \ne 0$ for the quadratic character mod an odd prime $p$. % v7:L013
NonVacuity.specialization proves that the challenge theorem specialises to an explicit statement about the quadratic characters. % v7:L014
Each of the four depends on the axioms [propext, Classical.choice, Quot.sound]. % v7:L011 v7:L012 v7:L013 v7:L014
\footnote{The lemmas are in lean-checks/NonVacuity.lean, LF sha256 b9fae4da5517a97da50d3fdb5356329f03f00b7ab935a448a999b2c59b24d9f9. % v7:L015
They were compiled as checker via sudo -u checker, with no systemd-run unit and no address-family filter. % v7:L138
}

\subsection{The three 7/8 challenges}

The other three challenges encode the main theorem of the 7/8 paper, ``The Quasi-Riemann Hypothesis'' (\texttt{\textbackslash date\{30 September 2026\}}, author OpenAI), whose paper.pdf has sha256 8fe93046f8cf5ef1ba5969c89addc02d76311adc4ee907509ff9cd96f7ec99e7 at both adc7f12 and fd4aeeb. % v7:L144 v7:L002 v7:L004 v7:L006 v8:L150
Lines 107-110 of its build/paper.tex read: % v7:L144
%% QUOTE external/openai-math/preprints/The-Quasi-Riemann-Hypothesis-September-30-2026/build/paper.tex 107-110 sha=v7:L144
\begin{paperquote}{}
Every finite-order Hecke $L$-function over
$F=\mathbb Q(\sqrt{-3})$ has no zero in $\Re s>7/8$.
The same holds for every Dirichlet $L$-function, including $\zeta(s)$.
A pole at $s=1$ for a principal character is allowed.
\end{paperquote}

Listing~\ref{lst:78-challenges} in Appendix~\ref{app:78-listings} quotes the three challenge statements. % v7:L001 v7:L003 v7:L005

By reading each Lean statement against the paper's text, all three match the paper statement; the QuasiRiemannHypothesis finding carries the qualification ``except s = 1, covered by Mathlib's junk value''. % v7:L002 v7:L004 v7:L006
The HeckeSevenEighths finding rests on a ruling about the L-function row and on one INFERRED identification. % v7:L006
The Lean sub-claims of these readings (non-vacuity, series identities, holomorphy) come from a separate statement-check build. % v7:L002 v7:L009
All 34 \texttt{\#print axioms} outputs of that build are exactly [propext, Classical.choice, Quot.sound]. % v7:L009
That build ran as checker via sudo -u checker, with no systemd-run unit and no address-family filter. % v7:L123
\footnote{The four statement audits were also posted publicly on grwtsk/openai-math \#29, at 2026-10-20T16:56:00Z. % v7:L010
}

\begin{table}[tbp]
\caption{The four family-003 challenge statements. File sha256 at openai/math fd4aeeb; finding and reading record as in the ledger. Every finding is INFERRED, by reading the Lean statement against the paper's text.}\label{tab:faithfulness} % v7:L001 v7:L002 v7:L003 v7:L004 v7:L005 v7:L006 v7:L007 v7:L008
\small
\begin{tabularx}{\textwidth}{@{}l>{\raggedright\arraybackslash}X>{\raggedright\arraybackslash}p{3.4cm}@{}}
\toprule
challenge & finding & reading \\ % no-fact: table header
\midrule
QuasiRiemannHypothesis & MATCH (except s = 1, covered by Mathlib's junk value) & Read by J. Reilly, 2026-10-27: MATCH \\ % v7:L002
\multicolumn{3}{@{}l}{\scriptsize\texttt{065f8c9a01d28db78c8b1bfc5083b535082a2d2aac2230d563f98b8caa812cc5}} \\ % v7:L001
DirichletSevenEighths & MATCH & Read by J. Reilly, 2026-10-27: MATCH \\ % v7:L004
\multicolumn{3}{@{}l}{\scriptsize\texttt{8fb13ad977bca104591fcbc6ef01a58496e7385053c9b1445f95749afc0eff54}} \\ % v7:L003
HeckeSevenEighths & MATCH (L-function row by Joe's ruling 2026-10-20; one identification INFERRED) & Read by J. Reilly, 2026-10-27: MATCH \\ % v7:L006
\multicolumn{3}{@{}l}{\scriptsize\texttt{cb5e404fa65a502c0a9006b9ad45c28914feb79e4b7b1ffcf05e4af82ed0cfc2}} \\ % v7:L005
SiegelZeros & MATCH & Read by J. Reilly, 2026-10-27: MATCH \\ % v7:L008
\multicolumn{3}{@{}l}{\scriptsize\texttt{bedaac0248a641910f79a380d9702d82c6bfbbfd3760a69d27c72c5d8fb136ce}} \\ % v7:L007
\bottomrule
\end{tabularx}
\end{table}

\subsection{Reasoning behind the readings}

Tables~\ref{tab:points-SiegelZeros}--\ref{tab:points-HeckeSevenEighths} list the reasoning points behind each challenge reading, one ledger row per point. % v8:L152 v8:L181
Each point is INFERRED, by reading the Lean statement against the paper's text, and none is yet marked by a second reader. % v8:L152 v8:L181

\begin{table}[!htbp]
\caption{Reasoning points behind the SiegelZeros reading (v7:L008). INFERRED; not yet marked by a second reader.}\label{tab:points-SiegelZeros} % v8:L152 v8:L153 v8:L154 v8:L155 v8:L156 v8:L157 v8:L158 v8:L159 v8:L160
\scriptsize\setlength{\tabcolsep}{3pt}\renewcommand{\arraystretch}{0.95}
\begin{tabularx}{\textwidth}{@{}>{\raggedright\arraybackslash}Xl@{}}
\toprule
point & finding \\ % no-fact: table header
\midrule
quantifier order: the constant c precedes q, chi and beta, so c is absolute & MATCH \\ % v8:L152
hypotheses: primitive, nonprincipal, real, q >= 3, beta in (0,1); none added, none missing & MATCH \\ % v8:L153
conclusion: c <= (1 - beta) log q, non-strict as in (1.1) & MATCH \\ % v8:L154
Mathlib meaning: IsPrimitive is conductor = level, and conductor is the least level through which chi factors & MATCH \\ % v8:L155
Mathlib meaning: (chi a).im = 0 for all a means real-valued; values are 0 or roots of unity, so 0 or +-1 & MATCH \\ % v8:L156
Mathlib meaning: DirichletCharacter.LFunction is ZMod.LFunction, the continuation; it equals the Dirichlet series for Re s > 1 and is entire for chi != 1 & MATCH \\ % v8:L157
edge cases: q = 4 (odd) and q = 8 (one even, one odd) are covered & MATCH \\ % v8:L158
edge cases: no hypothesis mentions parity, so both parities are included & MATCH \\ % v8:L159
edge cases: nothing excludes non-prime or even conductors & MATCH \\ % v8:L160
\bottomrule
\end{tabularx}
\end{table}

\begin{table}[!htbp]
\caption{Reasoning points behind the QuasiRiemannHypothesis reading (v7:L002). INFERRED; not yet marked by a second reader.}\label{tab:points-QuasiRiemannHypothesis} % v8:L161 v8:L162 v8:L163 v8:L164 v8:L165
\scriptsize\setlength{\tabcolsep}{3pt}\renewcommand{\arraystretch}{0.95}
\begin{tabularx}{\textwidth}{@{}>{\raggedright\arraybackslash}Xl@{}}
\toprule
point & finding \\ % no-fact: table header
\midrule
quantifier order: 7/8 is a fixed literal & MATCH \\ % v8:L161
hypotheses: only Re s > 7/8; no pole hypothesis & MATCH \\ % v8:L162
conclusion: != 0 on the open half-plane Re s > 7/8; the paper excludes the boundary & MATCH \\ % v8:L163
Mathlib meaning: riemannZeta is hurwitzZetaEven 0, the continuation, with a junk value at 1 & MATCH \\ % v8:L164
edge case s = 1: Mathlib's riemannZeta 1 is a nonzero junk value, so the statement asserts a true fact about it and nothing about the pole, which the paper allows & MATCH \\ % v8:L165
\bottomrule
\end{tabularx}
\end{table}

\begin{table}[!htbp]
\caption{Reasoning points behind the DirichletSevenEighths reading (v7:L004). INFERRED; not yet marked by a second reader.}\label{tab:points-DirichletSevenEighths} % v8:L166 v8:L167 v8:L168 v8:L169 v8:L170 v8:L171
\scriptsize\setlength{\tabcolsep}{3pt}\renewcommand{\arraystretch}{0.95}
\begin{tabularx}{\textwidth}{@{}>{\raggedright\arraybackslash}Xl@{}}
\toprule
point & finding \\ % no-fact: table header
\midrule
quantifier order: fixed 7/8, uniform over every q >= 1, chi and s & MATCH \\ % v8:L166
hypotheses: any modulus q >= 1 and any character, imprimitive ones and q = 1 included, as in the paper's definition & MATCH \\ % v8:L167
hypotheses: hpole excludes s = 1 only for the principal character mod q & MATCH \\ % v8:L168
conclusion: != 0, strict half-plane & MATCH \\ % v8:L169
Mathlib meaning: LFunction is the continuation; for q = 1 it is riemannZeta; for chi != 1 it is entire, so its value at s = 1 is genuine & MATCH \\ % v8:L170
edge cases: imprimitive nonprincipal characters gain zeros only on Re s = 0; the principal pole at 1 is excluded by hpole & MATCH \\ % v8:L171
\bottomrule
\end{tabularx}
\end{table}

\begin{table}[!htbp]
\caption{Reasoning points behind the HeckeSevenEighths reading (v7:L006). INFERRED; not yet marked by a second reader.}\label{tab:points-HeckeSevenEighths} % v8:L172 v8:L173 v8:L174 v8:L175 v8:L176 v8:L177 v8:L178 v8:L179 v8:L180 v8:L181
\scriptsize\setlength{\tabcolsep}{3pt}\renewcommand{\arraystretch}{0.95}
\begin{tabularx}{\textwidth}{@{}>{\raggedright\arraybackslash}Xl@{}}
\toprule
point & finding \\ % no-fact: table header
\midrule
quantifier order: fixed 7/8, every Character & MATCH \\ % v8:L172
hypotheses: Character (nonzero modulus, MulChar on O/m trivial on units, zero on non-units) is exactly a ray class character mod m extended by zero, since O is a PID with 6 units & MATCH \\ % v8:L173
hypotheses: imprimitive characters are included on both sides & MATCH \\ % v8:L174
hypotheses: the auxiliary period does not matter (HeckeHolo.period\_indep) & MATCH \\ % v8:L175
hypotheses: hpole excludes s = 1 for the principal character only & MATCH \\ % v8:L176
conclusion: != 0, strict half-plane & MATCH \\ % v8:L177
definitions: the lattice series for Re s > 1, holomorphy off \{0, 1\} and the norm formula fix the function in 7/8 < Re s <= 1 by the identity theorem & MATCH \\ % v8:L178
identification: each nonzero ideal has exactly 6 generators; chi and the norm are unit-invariant; so (1/6) times the sum over alpha equals the sum over ideals (absolute convergence) & MATCH \\ % v8:L179
the Dedekind factorization is not needed for this match (checked numerically only) & consistent \\ % v8:L180
edge case s = 0 (outside the half-plane): a Gamma(0) = 0 junk value, no effect on the statement & MATCH \\ % v8:L181
\bottomrule
\end{tabularx}
\end{table}
